% TOP_PROBLEM: {"id": 3072, "title": "Edge-antipodal colorings of cubes", "category_id": 2, "category": {"id": 2, "name": "combinatorics", "display_name": "Combinatorics", "description": "Counting problems, graph theory, discrete structures.", "slug": "combinatorics", "order_index": 2, "created_at": "2026-07-31T15:26:25.671Z"}, "status": "open", "problem_number": "OPG-2359", "set_id": 12}
% FIRST_POSTED: 2026-10-01
% ENTRY_KIND: statement_only
% UnsolvedMath ID 3072; source catalogue code OPG-2359
% !TeX program = lualatex
% Definition and sources only; this file is not an LLM solution attempt.
\documentclass[11pt,a4paper]{article}
\usepackage[margin=25mm]{geometry}
\usepackage{fontspec}
\setmainfont{lmroman10-regular.otf}[BoldFont=lmroman10-bold.otf,
 ItalicFont=lmroman10-italic.otf,BoldItalicFont=lmroman10-bolditalic.otf]
\usepackage{amsmath,amssymb,mathtools}
\usepackage{unicode-math}
\setmathfont{latinmodern-math.otf}
\AtBeginDocument{\let\setminus\smallsetminus}
\usepackage{xurl,hyperref}
\hypersetup{colorlinks=true,urlcolor=blue}
\setcounter{secnumdepth}{0}
\setlength{\parindent}{0pt}
\setlength{\parskip}{5pt}
\setlength{\emergencystretch}{3em}
\begin{document}
\section*{Edge-antipodal colorings of cubes}\label{3072}
{\small UnsolvedMath ID 3072 · OPG-2359 · Combinatorics · OpenGarden}
\subsection{Definitions and mathematical statement}
For $d\ge2$, let $Q_d$ be the graph on $\{0,1\}^d$ whose edges join
strings differing in exactly one coordinate. For a vertex
$x=(x_1,\ldots,x_d)$, define its antipode by
$\bar x=(1-x_1,\ldots,1-x_d)$. The antipodal edge to an unordered edge
$e=\{x,y\}$ is $\bar e=\{\bar x,\bar y\}$, which is again an edge.

An edge-antipodal two-colouring is a map
$\varphi:E(Q_d)\to\{0,1\}$ satisfying
$\varphi(\bar e)=1-\varphi(e)$ for every edge $e$.
No proper-colouring condition is imposed at a vertex.
The conjecture asserts that for every dimension and every such colouring
there are a vertex $x$ and a path from $x$ to $\bar x$ all of whose edges
have the same colour.

More explicitly, one seeks distinct vertices
$v_0,\ldots,v_t$ with $v_0=x$, $v_t=\bar x$, adjacent consecutive
vertices, and one $b\in\{0,1\}$ such that
$\varphi(\{v_{i-1},v_i\})=b$ for all $1\le i\le t$.
There is no requirement that $t=d$, that the path be shortest, or that
it change each coordinate just once.

Equivalently, some connected component of one of the two monochromatic
spanning subgraphs must contain an antipodal vertex pair. The pair is
chosen after the colouring and need not be a prescribed pair. Antipodality
is taken in the full $d$-cube; restricting the colouring to a lower-dimensional
face need not preserve this hypothesis.
\subsection{Short English statement}
Must a cube edge-colouring that reverses colour on antipodal edges contain a monochromatic path between some antipodal vertices?
\subsection{Sources}
\begin{itemize}
\item[S1] ULAM.AI and the UnsolvedMath Contributors, supplied problems.json, record 3072 (OPG-2359), OpenGarden collection. Catalogue identity and source transcription; CC BY 4.0.
\url{https://huggingface.co/datasets/ulamai/UnsolvedMath}.
\item[S2] Open Problem Garden, Edge-antipodal colorings of cubes. Original problem and discussion; the mathematical formulation below is expanded from the supplied source record.
\url{http://www.openproblemgarden.org/op/edge_antipodal_colorings_of_cubes}.
\end{itemize}
\end{document}
